\section{Angular momentum and the role of symmetry}
\label{r:sec:angular}

The distinction between the leading vortex and the completed velocity
matters to the comparison. The background in OpenAI's construction is
axisymmetric. Its wave potentials contain factors
\(e^{ikm\Phi}\), with \(\Phi=p\theta+\varphi\) and nonzero integer
angular frequency \(kmp\)
\cite[(6.27), (7.3), (7.38)]{openai-paper}. Taking their curls retains
these harmonics. They enter the realized and localized velocity through
\cite[(9.21), (10.4)]{openai-paper}; the final spatial cutoff does not
average them over the angle. The complete field is nonaxisymmetric.

One can see exactly why that distinction affects a force bound. Define
angular momentum per unit mass about the axis by
\(\Gamma=ru_\theta\). The angular component of \eqref{r:eq:ns}, including
the derivatives of the cylindrical basis vectors, gives for \(r>0\)
\begin{equation}
 \begin{aligned}
 \left(\partial_t+u_r\partial_r
       +\frac{u_\theta}{r}\partial_\theta+u_z\partial_z\right)\Gamma
 &=\nu\left(\partial_{rr}-\frac1r\partial_r
       +\frac1{r^2}\partial_{\theta\theta}+\partial_{zz}\right)\Gamma\\
 &\quad-\partial_\theta p
       +\frac{2\nu}{r}\partial_\theta u_r+rf_\theta.
 \end{aligned}
 \label{r:eq:angular-full}
\end{equation}
Indeed, multiplying the angular momentum equation for \(u_\theta\) by
\(r\) combines its geometric term \(u_ru_\theta\) with the material
derivative of \(r\). The vector Laplacian contributes
\(2\nu r^{-1}\partial_\theta u_r\). Pressure supplies the torque
\(-\partial_\theta p\). These terms describe internal momentum transfer;
they are present even when the prescribed force is small.

\begin{proposition}[A force bound for axisymmetric angular momentum]
\label{r:prop:angular-max}
Suppose \(u,p,f\) are axisymmetric, solve \eqref{r:eq:ns} from rest, and
\(u\) is bounded on every \(\R^3\times[0,S]\), \(S<T\).
If \(G(t):=\int_0^t\norm{rf_\theta(s)}_\infty\dd s\) is finite through
\(T\), then
\begin{equation}
 \norm{ru_\theta(t)}_\infty\le G(t),\qquad 0\le t<T.
 \label{r:eq:angular-max}
\end{equation}
There is no source-amplitude restriction in this bound.
\end{proposition}

\begin{proof}
Under these symmetry assumptions the angular derivatives in
\eqref{r:eq:angular-full} vanish. The resulting scalar operator is
\[
 L=\partial_t+u_r\partial_r+u_z\partial_z
       -\nu(\partial_{rr}-r^{-1}\partial_r+\partial_{zz}).
\]
Fix \(S<T\) and let \(M_S=\sup_{[0,S]\times\R^3}|u|\).
For \(\Psi=1+r^2+z^2\) and \(c=M_S+2\nu+1\), direct differentiation
gives
\[
 L(e^{ct}\Psi)
 =e^{ct}\bigl(c\Psi+2ru_r+2zu_z-2\nu\bigr)>0.
\]
Here \(2M_S\sqrt{r^2+z^2}\le M_S\Psi\).
Compare \(\pm\Gamma\) with \(G(t)+\epsilon e^{ct}\Psi\).
The comparison functions dominate initially. At the axis,
\(\Gamma=0\) by Cartesian smoothness; at spatial infinity their
quadratic growth dominates \(|\Gamma|\le M_Sr\).
Apply the scalar maximum principle on bounded cylinders away from the
axis, using the integral of the maximum of \(|rf_\theta|\) on each
cylinder as the temporal barrier. That integral is at most \(G\).
For sufficiently large outer radius the quadratic barrier dominates the
outer boundary values; at a sufficiently small excluded axial radius,
\(|\Gamma|\le M_Sr\) is dominated by the positive barrier.
Exhausting the domain gives
\(\pm\Gamma\le G+\epsilon e^{ct}\Psi\) on \([0,S]\).
Let \(\epsilon\downarrow0\) and then use the arbitrariness of \(S<T\).
\end{proof}

Smooth rapidly decreasing forcing satisfies the hypothesis on \(G\).
This estimate would exclude an axisymmetric realization of the exterior
growth calculated in Section~\ref{r:sec:exact-cancellation}, since there \(ru_\theta\) grows like
\((1-t)^{-h}\). It does not assert global regularity for every
axisymmetric velocity: a bound on \(ru_\theta\) is not a uniform bound
on \(u_\theta\) near the axis.

Nor does angular averaging turn the completed nonaxisymmetric field into
a solution of the axisymmetric equation. Write a bar for angular average
and a prime for the fluctuation about that average. Averaging the
conservative form of \eqref{r:eq:angular-full} gives exactly
\begin{equation}
 \begin{aligned}
 \partial_t\bar\Gamma+\bar u_r\partial_r\bar\Gamma
       +\bar u_z\partial_z\bar\Gamma
 &=\nu(\partial_{rr}-r^{-1}\partial_r+\partial_{zz})\bar\Gamma
       +r\bar f_\theta\\
 &\quad-\frac1r\partial_r
       \left(r\overline{u_r'\Gamma'}\right)
       -\partial_z\overline{u_z'\Gamma'}.
 \end{aligned}
 \label{r:eq:angular-mean}
\end{equation}
The pressure torque and angular viscous derivative average to zero, but
the products in the transport flux do not. Two zero-mean harmonics can
have a nonzero averaged product. OpenAI's covariance construction uses
precisely these products to supply momentum flux
\cite[(7.23), (9.11)]{openai-paper}. A force-only estimate for
\(\bar\Gamma\) would have to control the two fluctuation fluxes in
\eqref{r:eq:angular-mean}. Dropping them would change the equation being
compared.
